\documentclass[12pt]{article}
\usepackage{amssymb,amsmath}
\usepackage{anyfontsize}
\thispagestyle{empty}
\usepackage{tikz}
\usetikzlibrary{calc}
\usepackage{setspace}
\usepackage{graphicx}
\usepackage[margin=0.65in]{geometry}
\usepackage{float}
\usepackage[letterspace=150]{microtype}


\begin{document}

%Border

\begin{tikzpicture}[overlay,remember picture]
\draw [line width=3pt,rounded corners=0pt,]
    ($ (current page.north west) + (1.25cm,-1.25cm) $)
    rectangle
    ($ (current page.south east) + (-1.25cm,1.25cm) $);       
\end{tikzpicture}
\begin{center}\includegraphics[scale=.35]{RUMAlogo.png}\\
presents... \\
\begin{spacing}{1}
{\fontsize{40}{44}\selectfont  \textsc{
Connection Games. Cool Stuff.}} \end{spacing}
 

~~\\
\begin{spacing}{1.5}
{\fontsize{24}{24} \selectfont A Lecture by Professor David Molnar}  \end{spacing}
\large Dept. of Mathematics, Rutgers University \\~~\\
~~\\

\normalsize
\textsc{Abstract:}

\Large 
This is, at its root, a talk about finding Mathematics where you
look. Since I like to play games, I tend to find Math there. (But
this is not a talk about game theory.)
I will tell you about some games -- including Hex, Y, Star Y, and
Atoll -- known as {\it connection games} (because the goal is to
connect things). So, even if you don't follow all the theory, you
can leave with some new games to play. I will also prove, for at
least one of these games, that draws are not possible. This has
strategic implications. These proofs also make use, in one form or
another, of the staggeringly popular Sperner's Lemma.\\

\vspace{10mm} 
\Huge   \textsc{Tuesday\\April 30, 2019 \\Hill 705 at 7:00 pm} \\~\\
\large
\emph{*Pizza and refreshments will be served}
\end{center}



\end{document}



